\section{Lean connection to the actual Riemann hypothesis}
\label{sec:leanxi}
The Lean formalization now reaches the \emph{genuine}
completed zeta function rather than just a symbolic
Gaussian model. Mathlib's
\texttt{completedRiemannZeta\(\_0\)}
represents the entire pole-subtracted completion
$\Lambda_0(s)$, satisfying
\[
\Lambda_0(1-s)=\Lambda_0(s),\qquad
\Lambda_0(s)=\Lambda(s)+\frac1s+\frac1{1-s}
\]
off the poles of $\Lambda(s)=\pi^{-s/2}\Gamma(s/2)\zeta(s)$.
Thus the identity
\begin{equation}
\xi(s)=\frac{1+s(s-1)\Lambda_0(s)}{2}
\label{eq:xi_mathi}
\end{equation}
defines an entire function with the correct extension
at $s=0$ and $s=1$. Its critical-line reparameterization is
\[
\widehat\Xi(z)=\xi\!\left(\frac12+\frac{iz}{2}\right).
\]
The Lean modules \texttt{XiKernel/EntireXi.lean} and
\texttt{XiKernel/AxiomAudit.lean} formally check
the $\xi$-functional equation, $\xi(0)=\xi(1)=1/2$,
complex differentiability, and the evenness and
entireness of $\widehat\Xi$.
The exact convention used earlier is
$F(z)=\widehat\Xi(z)/4$.

The proof target is to establish the
\emph{real-zero property}
\[
\widehat\Xi(z)=0\quad\Longrightarrow\quad
\operatorname{Im}z=0.
\]
To turn this into Mathlib's native
\texttt{RiemannHypothesis}, the formalization
must additionally relate zeros of $\xi$ to the
nontrivial zeros of $\zeta$ using the completed
zeta representation, with the trivial zeros
and poles treated explicitly. That zero
correspondence is a separate outstanding
proof obligation. The present Lean results
do not establish the real-zero property.
